\begin{abstract}
Learned particle simulators usually inherit their interaction graphs from geometry, although extra messages can help some predictions and harm others. We develop constructive adaptive interaction graphs that preserve native edges and append an exact budget of local pairs. Cached residual scores allocate that budget without a second scoring pass after initialization. We pair this controller with graph-exposure training, separating learning to use extra messages from choosing where to place them. Exploratory three-seed studies on Goop, WaterDrop, and Sand compare both training arms and declared policies on common observed histories and full autonomous rollouts; WaterDrop uses paired continuations. Exposure improves fixed-random rollout error in every Goop and WaterDrop seed, reducing Goop mean MSE from 0.301 to 0.129. Risk-guided placement does not improve consistently: on Sand, exposure narrows the observed risk deficit while worsening cached-risk rollout error. We formalize the distinction between residual risk and interaction benefit, and separate expansion cost from changes in rollout geometry. The results identify training support, placement, and autonomous feedback as distinct requirements for effective adaptive graphs.
\end{abstract}
